\section{Studi Kasus: Pemodelan Aturan E-Commerce}

Logika predikat tidak hanya berguna dalam konteks akademik, tetapi juga diterapkan secara luas dalam sistem informasi. Studi kasus ini menunjukkan bagaimana aturan bisnis (\textit{business rules}) pada platform e-commerce dapat dimodelkan dan divalidasi menggunakan logika predikat orde pertama \cite{rosen2019}.

\subsection{Skenario Bisnis}

Sebuah platform e-commerce memiliki aturan-aturan berikut:
\begin{enumerate}
    \item Pengguna yang telah login dan terverifikasi dapat menyimpan item ke keranjang belanja.
    \item Pengguna yang memiliki keranjang berisi setidaknya satu item dan saldo mencukupi dapat melakukan checkout.
    \item Pengguna premium mendapatkan diskon otomatis pada setiap pembelian.
\end{enumerate}

\subsection{Pemodelan Formal}

Fakta ($\Delta$):
\begin{align*}
\Delta = \{ & Login(UserA), \; Terverifikasi(UserA), \\
            & Premium(UserA), \; Saldo(UserA, 500000), \\
            & HargaItem(Laptop, 450000), \; Simpan(UserA, Laptop, Keranjang) \}
\end{align*}
Fakta $Simpan(UserA, Laptop, Keranjang)$ merepresentasikan aksi pengguna yang telah menyimpan Laptop ke keranjang; aksi ini dicatat sebagai fakta, bukan diturunkan dari aturan.

Aturan ($\Gamma$):
\begin{align*}
R_1&: \forall x\, (Login(x) \wedge Terverifikasi(x) \rightarrow BolehSimpan(x)) \\
R_2&: \forall x, i, s, h\, (BolehSimpan(x) \wedge Simpan(x, i, Keranjang) \wedge Saldo(x, s) \wedge \\
    & \qquad\qquad HargaItem(i, h) \wedge (s \geq h) \rightarrow BolehCheckout(x, i)) \\
R_3&: \forall x, i\, (BolehCheckout(x, i) \wedge Premium(x) \rightarrow Diskon(x, i, 10\%))
\end{align*}

\subsection{Inferensi: Apakah UserA Berhak Diskon?}

\begin{contoh}
\textbf{Query}: $Diskon(UserA, Laptop, 10\%)$?

\textbf{Penalaran (Forward Chaining):}
\begin{enumerate}
    \item $Login(UserA) \wedge Terverifikasi(UserA)$ \hfill (fakta dari $\Delta$)
    \item $BolehSimpan(UserA)$ \hfill ($R_1$, Modus Ponens)
    \item $Simpan(UserA, Laptop, Keranjang)$ \hfill (fakta dari $\Delta$)
    \item $Saldo(UserA, 500000) \wedge HargaItem(Laptop, 450000) \wedge (500000 \geq 450000)$
    \hfill (fakta $\Delta$ + evaluasi)
    \item $BolehCheckout(UserA, Laptop)$ \hfill ($R_2$, Modus Ponens dari langkah 1--4)
    \item $BolehCheckout(UserA, Laptop) \wedge Premium(UserA)$ \hfill (langkah 5 + fakta $\Delta$)
    \item $Diskon(UserA, Laptop, 10\%)$ \hfill ($R_3$, Modus Ponens)
\end{enumerate}

Query terjawab: UserA berhak mendapat diskon 10\% untuk pembelian Laptop. $\blacksquare$
\end{contoh}

\subsection{Kasus Negatif: UserB Tidak Terverifikasi}

\begin{contoh}
Misalkan $\Delta' = \{Login(UserB), \neg Terverifikasi(UserB)\}$.

\textbf{Query}: $BolehSimpan(UserB)$?

Aturan $R_1$ mensyaratkan $Login(x) \wedge Terverifikasi(x)$. Karena $Terverifikasi(UserB)$ tidak ada dalam $\Delta'$ (atau secara eksplisit bernilai salah), anteseden $R_1$ tidak terpenuhi untuk $x = UserB$. Maka $BolehSimpan(UserB)$ \textbf{tidak dapat disimpulkan}. Rantai inferensi berhenti, dan UserB tidak dapat melanjutkan ke tahap checkout.
\end{contoh}

\begin{catatan}
Contoh kasus negatif ini menunjukkan bahwa basis pengetahuan menolak menyimpulkan hak akses ketika prekondisi (dalam hal ini verifikasi pengguna) tidak terpenuhi. Pola ini analog dengan mekanisme otorisasi (\textit{authorization}) pada sistem perangkat lunak.
\end{catatan}
